\section{多源能源供给模型}
\label{sec:renewables-load}

\subsection{模型边界与公共测点}

多源供给模块把同场风、光、潮环境数据、设备功率曲线和运行状态映射为公共汇集点的逐源
可用功率。该模块只发布源侧物理边界，不决定外送、储能、制氢、算力和海洋用能之间的
功率分配。图~\ref{fig:source-model}给出环境输入、单源物理映射、状态与损耗处理以及
统一测点之间的关系。
\begin{figure}[H]
\centering
\resizebox{0.96\textwidth}{!}{\begin{tikzpicture}[
  font=\small,>={Stealth[length=2mm]},
  box/.style={draw=green!55!black,fill=green!8,rounded corners=2pt,
    minimum width=27mm,minimum height=10mm,align=center},
  process/.style={draw=blue!60!black,fill=blue!7,rounded corners=2pt,
    minimum width=31mm,minimum height=12mm,align=center},
  output/.style={draw=teal!65!black,fill=teal!8,rounded corners=2pt,
    minimum width=31mm,minimum height=13mm,align=center},
  flow/.style={->,thick,draw=gray!80},
  info/.style={->,thick,dashed,draw=blue!65}
]
\node[box] (met) at (0,2.2) {风速/波浪/姿态};
\node[box] (solar) at (0,0.7) {辐照/海温/姿态};
\node[box] (tidal) at (0,-0.8) {流速/潮向/海况};
\node[process] (wind) at (3.5,2.2) {风机功率曲线\\与状态机};
\node[process] (pv) at (3.5,0.7) {组件辐照--温度\\与逆变器边界};
\node[process] (tc) at (3.5,-0.8) {涨/落潮曲线\\与换向状态};
\node[process] (loss) at (7.1,0.7) {可用率、降额、辅机\\与集电损耗归属};
\node[output] (poi) at (11,0.7) {MP-03-SOURCE-POI\\可用/实际/弃电功率};
\node[box] (dispatch) at (7.1,-1.45) {4.9调度请求};

\draw[flow] (met) -- (wind);
\draw[flow] (solar) -- (pv);
\draw[flow] (tidal) -- (tc);
\draw[flow] (wind) -- (loss);
\draw[flow] (pv) -- (loss);
\draw[flow] (tc) -- (loss);
\draw[flow] (loss) -- (poi);
\draw[info] (dispatch) -- (loss);
\end{tikzpicture}
}
\caption{多源能源供给子模型结构}
\label{fig:source-model}
\sourcetype{项目内部子模型结构图。}
\end{figure}

设备端毛功率到公共测点的换算为
\begin{equation}
P_{i,t}^{av}=\max\!\left(0,
P_{i,t}^{av,gross}-P_{i,t}^{loss,col}\right),
\qquad i\in\mathcal S,
\label{eq:source-poi}
\end{equation}
其中$P_{i,t}^{loss,col}$为设备端至公共汇集点的集电损耗。式~\eqref{eq:source-poi}
属于本项目计量点定义；损耗须由设计计算、效率曲线或计量对测校准，未校准时为
\assumptiontag。设备辅机功率作为独立母线消耗提交，不隐藏在发电注入中。

\subsection{漂浮式风电}

调度层优先采用拟选机组的OEM或型式试验电功率曲线。IEC~61400-12系列规定功率性能
测量口径，IEC~61400-3-2和GB/Z~44047规定浮式海上风机的设计边界
\cite{iec61400-12,iec61400-3-2,gbz44047}；这些标准不直接提供本项目机组功率曲线。

若离线浮式气动--水动--结构模型已提供平台运动，可定义沿转子轴的等效相对风速代理：
\begin{equation}
v_{rel,t}=\max\!\left[0,
v_{hub,t}-\dot x_{surge,t}-H_{hub}\dot\theta_{pitch,t}\right].
\label{eq:wind-relative}
\end{equation}
式~\eqref{eq:wind-relative}的坐标方向和高阶耦合项仍须与平台坐标系及OpenFAST离线结果
核对，现阶段标记为\verifytag。机组可用毛功率由经验证的电功率曲线插值得到：
\begin{equation}
P_{w,t}^{curve}=f_{P-v}(v_{rel,t}),\qquad
0\le P_{w,t}^{av,gross}\le N_wP_w^r.
\label{eq:wind-curve}
\end{equation}
若输入曲线已经是机端电功率，不再重复乘发电机或变流器效率。场级尾流、设备可用率、
环境降额、辅机与集电损耗分别记录，不合并成不可解释的固定效率。

风机状态机至少区分正常、资源不足、高风速停机、复归等待、故障/检修和降额。正常受控
条件下的实际利用功率满足
\begin{equation}
0\le P_{w,t}^{use}\le P_{w,t}^{av},
\label{eq:wind-available}
\end{equation}
\begin{equation}
-R_w^{dn}\Delta t_t\le P_{w,t}^{use}-P_{w,t-1}^{use}
\le R_w^{up}\Delta t_t.
\label{eq:wind-ramp}
\end{equation}
资源骤降、保护停机或故障必须优先服从实时可用边界，不能因下降爬坡约束而制造高于
可用功率的输出。切出、复归和延时参数须来自OEM控制规范或认证试验，未取得时为
\assumptiontag。

\subsection{漂浮式光伏}

平台滚转、纵摇和艏摇改变组件法向量$\bm n_t$。考虑直射、天空散射、海面反射与双面
背面辐照的调度代理可写为
\begin{equation}
\begin{aligned}
G_{POA,t}={}&G_{DNI,t}\max(0,\bm s_t\!\cdot\!\bm n_t)
+G_{DHI,t}\frac{1+n_{z,t}}{2}\\
&+\rho_gG_{GHI,t}\frac{1-n_{z,t}}{2}
+\varphi_bG_{rear,t}.
\end{aligned}
\label{eq:pv-poa}
\end{equation}
式~\eqref{eq:pv-poa}采用各向同性天空和海面反射近似；镜面反射、遮挡、波面统计及姿态
平均方法尚需实测校准，标记为\verifytag。优先使用IEC~61853系列试验形成的辐照--温度
功率矩阵\cite{iec61853}；缺少矩阵时，可暂用厂家温度系数代理
\begin{equation}
P_{dc,t}=P_{dc}^{r}\frac{G_{POA,t}}{G_{STC}}
\left[1+\gamma_P(T_{m,t}-T_{STC})\right]
k_{soil,t}k_{deg,t}k_{mis,t}.
\label{eq:pv-dc}
\end{equation}
交流侧功率及逆变器简化边界为
\begin{equation}
P_{pv,t}^{ac}=\min\!\left(P_{ac}^{r},\eta_{inv}(l_t)P_{dc,t}\right),
\qquad (P_{pv,t}^{ac})^2+Q_{pv,t}^2\le(S_{pv}^{r})^2.
\label{eq:pv-ac}
\end{equation}
工程应用应以厂家P--Q能力图替代简单视在功率圆。海上盐雾污损、腐蚀、阵列姿态、波致
失配、停机和复归阈值均需OEM或海上实证校准；NB/T~11814可作为设计边界参考
\cite{nbt11814}，不能替代设备认证数据。

\subsection{潮流能预留模型}

潮流速度保留涨落潮方向。考虑平台轴向运动和阵列尾流后的等效流速为
\begin{equation}
u_{rel,t}=(u_t-u_{p,t})k_{wake,t},
\label{eq:tidal-relative}
\end{equation}
其中$k_{wake,t}$须由CFD、阵列试验或现场辨识获得，未校准时为\assumptiontag。设备级
功率应优先采用GB/T~41342对应的现场功率特性测试或厂家电功率曲线
\cite{gbt41342}，并对涨、落潮分别插值：
\begin{equation}
P_{tc,t}^{curve}=
\begin{cases}
f_{flood}(|u_{rel,t}|), & u_{rel,t}>u_{db},\\
f_{ebb}(|u_{rel,t}|), & u_{rel,t}<-u_{db},\\
0, & |u_{rel,t}|\le u_{db}.
\end{cases}
\label{eq:tidal-curve}
\end{equation}
若曲线已为电功率，不再重复乘$C_p$及传动链效率。IEC~TS~62600-200/201可用于功率
性能和资源评估口径\cite{iec62600}；换向、切入/切出流速、波高停机和复归时间仍须由
OEM或示范项目数据确定。

\subsection{多源聚合、不确定性与状态输出}

各源的可用、已用与弃电必须分开记录：
\begin{equation}
0\le P_{i,t}^{use}\le P_{i,t}^{av},\qquad
P_{i,t}^{curt}=P_{i,t}^{av}-P_{i,t}^{use}.
\label{eq:source-use-curt}
\end{equation}
在统一测点的聚合量为
\begin{equation}
P_{src,t}^{x}=\sum_{i\in\mathcal S}P_{i,t}^{x},
\qquad x\in\{av,use,curt\}.
\label{eq:source-aggregate}
\end{equation}
调度弃电、资源不足、环境停机、故障检修、集电损耗和辅机耗电应分别统计，不得全部记作
“弃风弃光”。理论上调余量$P_{src,t}^{av}-P_{src,t}^{use}$还需结合持续时间、预测误差、
通信、爬坡和设备状态后，才能作为可靠备用。

概率场景$k\in\mathcal K$写为
\begin{equation}
\bm P_{t,k}^{scen}=
[P_{w,t,k},P_{pv,t,k},P_{tc,t,k}]^{\mathsf T},\qquad
\sum_{k\in\mathcal K}\pi_k=1,\quad \pi_k\ge0.
\label{eq:joint-source-scenario}
\end{equation}
场景必须保持源间相关性、时间连续性、昼夜周期、潮汐相位和共同极端事件。历史轨迹
重采样、联合概率预测或Copula等方法只有在完成独立滚动回测后方可用于工程结论；在核验
具体概率模型原文之前，其公式统一标记为\verifytag。

本模块最终向联合模型提供逐源可用功率、可用状态、爬坡、无功能力、质量标志和事件状态。
本节不呈现合成场景测试、互补性算例或验证结果；相关测试与图表统一放入第5章。
